\begin{abcliste}
\abc
Auf das Brett wirken die Gewichtskraft, die im Schwerpunkt S angreift, und die Kraft der Achse in D. Bezüglich D hat die Kraft der Achse keinen Hebelarm. Die Wirkungslinie der Gewichtskraft ist die Senkrechte durch S, und die geht durch D: Auch die Gewichtskraft hat keinen Hebelarm und kein Drehmoment. Das Brett bleibt \result{\text{in Ruhe}}.

S liegt unter D. Dreht man das Brett ein wenig, wandert S auf einem Kreisbogen um D zur Seite und nach oben, und die Gewichtskraft dreht es zurück: \result{\text{stabiles Gleichgewicht}}.

\abc
Ob ein Körper an einer Achse in Ruhe bleibt, hängt nicht davon ab, wie seine Kanten liegen, sondern allein davon, wo sein Schwerpunkt bezüglich der Achse liegt: senkrecht darunter, senkrecht darüber oder in der Achse. Ein Brett, dessen Achse nicht in der Mitte liegt, hängt in seiner stabilen Lage im Allgemeinen schief. Waagrecht hängt es nur, wenn D senkrecht über der Mitte liegt.

\abc
Das Brett pendelt um die gezeichnete Lage hin und her (ohne Reibung immer weiter, mit etwas Reibung immer schwächer) und kommt mit Reibung wieder in genau dieser schiefen Lage zur Ruhe.
\end{abcliste}
